\section{Surgery limits and the first-order variational problem}\label{sec:variational-geometry}

We apply the linear estimate to an explicit sequence of closed manifolds.
The surgery construction gives a recovery sequence, but does not by itself
determine the best first-order constant.  We first retain the geometric
information needed for that construction, and then identify the unrestricted
variational problem to which the constant belongs.

\subsection{An actual surgery class and its curvature budgets}

Let $X$ be a closed complex-hyperbolic surface, with its complex orientation,
and put $A=c_1^2(X)>0$.  Normalize its complex-hyperbolic metric $g_X$ by
$\Ric_{g_X}=-g_X$.  Its characteristic numbers and curvature integrals are
\begin{equation}\label{eq:vg-background}
 \begin{gathered}
 \chi(X)=\frac A3,\qquad \tau(X)=\frac A9,\qquad
 \Vol(g_X)=2\pi^2A,\\
 s_{g_X}=-4,\qquad W^-_{g_X}=0,
 \qquad \int_X|W^+_{g_X}|^2=\frac{4\pi^2A}{3}.
 \end{gathered}
\end{equation}
Here and below Ricci norms are tensor norms, whereas Weyl norms are
Hilbert--Schmidt norms of endomorphisms of two-forms.  The volume formula
follows from $[\omega_X]=-2\pi c_1(X)$ and
$\Vol(g_X)=\frac12[\omega_X]^2$.  The Weyl formula follows from the
K\"ahler identity $|W^+|^2=s^2/24$.

For integers $k,\ell\geq0$, define
\begin{equation}\label{eq:vg-surgery-data}
 \begin{gathered}
 M=X\#k\overline{\mathbb{CP}}^{\,2}\#\ell(S^1\times S^3),
 \qquad Q=k+4\ell,\qquad h=\frac QA,\\
 T_+=2\chi(M)+3\tau(M)=A-Q,\qquad
 T_-=2\chi(M)-3\tau(M)=\frac A3+5k-4\ell,\\
 U=\frac{3\min\{T_+,T_-\}}A.
 \end{gathered}
\end{equation}
Indeed, a blow-up increases $\chi$ by one and decreases $\tau$ by one,
whereas an $S^1\times S^3$ summand decreases $\chi$ by two and leaves
$\tau$ unchanged.  Thus
$\chi(M)=A/3+k-2\ell$ and $\tau(M)=A/9-k$, which give
\eqref{eq:vg-surgery-data}.

The gauge-theoretic input on these manifolds is an actual polarized
Seiberg--Witten class, not just a characteristic element of the intersection
form.  For each metric $g$, set $H=H_g^+$ and pull
$c=-K_X$ back to $M$.  Choose signs of the exceptional generators so that
\begin{equation}\label{eq:vg-sign-selection}
 E_i\cdot E_j=-\delta_{ij},\qquad c\cdot E_i=0,
 \qquad c_H^+\cdot E_i\leq0.
\end{equation}
The polarized blow-up and handle construction in
\cite[Proposition~3.2 and its proof]{LeBrun2001} supplies a
Spin$^c$ structure with determinant class
\begin{equation}\label{eq:vg-actual-class}
 a_H=c-\sum_{i=1}^k E_i.
\end{equation}
Its invariant is nonzero, with the $\ell$ degree-one loop insertions from
the handle summands.  Consequently its unperturbed monopole equations
are solvable for every $H$-adapted metric.  When $b^+=1$, this assertion
uses the chamber selected in that polarized construction.  It is not an
assertion of chamber-independent nonvanishing for an arbitrary class.
For the canonical class on $X$, its positive square keeps the canonical
small-perturbation chamber away from the wall; the sign choice in
\eqref{eq:vg-sign-selection} preserves that chamber on the connected sum.

The numerical consequences, with the structure now fixed, are
\begin{equation}\label{eq:vg-actual-budget}
 \begin{split}
 p:=(a_H^+)^2
   &=(c_H^+)^2-2\sum_i c_H^+\cdot E_i
                        +\left|\sum_iE_i^+\right|^2
     \geq(c_H^+)^2\geq A,\\
 a_H^2&=A-k\geq A-k-4\ell=T_+.
 \end{split}
\end{equation}
Since the Hodge star on two-forms is conformally invariant, $H$ and the
chosen Spin$^c$ structure remain fixed throughout the conformal class.
In particular this is the actual conformal supply required by
Theorem~\ref{thm:linear-window}.  A different initial metric is allowed
to select a different structure.  The handles contribute to $T_+$, but
do not decrease the square in \eqref{eq:vg-actual-budget}.

For use in this section, write
\begin{equation}\label{eq:vg-energies}
 \begin{gathered}
 V=\Vol(g),\qquad I=\int_Ms_g\,d\mu_g,\qquad
 S=\int_Ms_g^2\,d\mu_g,\\
 P=\int_M|\Ric_g^\circ|^2\,d\mu_g,\qquad
 R=\int_M|\Ric_g|^2\,d\mu_g=P+\frac S4,\qquad
 L_s=S-\frac{I^2}{V}=\int_M(s_g-I/V)^2\,d\mu_g.
 \end{gathered}
\end{equation}
The scalar monopole estimate is
\begin{equation}\label{eq:vg-scalar-budget}
 S\geq32\pi^2p\geq32\pi^2A.
\end{equation}
With the normalization $|\mathfrak q(\Phi)|^2=|\Phi|^4/8$, the
unperturbed equations and the integrated spinor formula give
\[
 32\pi^2p\leq H_\Phi:=\int_M|\Phi|^4,
 \qquad
 4\int_M|\nabla_{\mathcal A}\Phi|^2
       +\int_Ms|\Phi|^2+H_\Phi=0.
\]
The first inequality is the harmonic projection bound for the self-dual
determinant curvature.  Cauchy--Schwarz in the second gives
$H_\Phi\leq S^{1/2}H_\Phi^{1/2}$.  Since $p>0$, one has $H_\Phi>0$,
and \eqref{eq:vg-scalar-budget} follows.  In particular $R>0$ for every
metric under consideration.

Minimizing the quadratic in $1/\lambda$ yields
\begin{equation}\label{eq:vg-defect-formulas}
 \E(g)=\inf_{\lambda\ne0}
   \frac{\int_M|\Ric_g-\lambda g|^2\,d\mu_g}{\lambda^2V}
   =4-\frac{I^2}{VR}
   =\frac{4P+L_s}{R},\qquad
 \J(g)=\frac{4P}{R}\leq\E(g).
\end{equation}
Indeed, the expression to be minimized is
$4-2I/(\lambda V)+R/(\lambda^2V)$.  If $I=0$, the value four is
an infimum as $|\lambda|\to\infty$; the same formula therefore applies.

The two Chern--Weil identities read
\begin{equation}\label{eq:vg-chiral-identities}
 W_\pm:=\int_M|W_g^\pm|^2\,d\mu_g
       =2\pi^2T_\pm-\frac{R}{12}+\frac P3.
\end{equation}
Their nonnegativity and \eqref{eq:vg-scalar-budget} imply, for every metric,
\begin{equation}\label{eq:vg-two-sided-R}
 8\pi^2A+P\leq R\leq8\pi^2AU+4P.
\end{equation}
Equivalently, with $e=R/(8\pi^2A)$, the upper bound is
\begin{equation}\label{eq:vg-topological-ceiling}
 e(1-\J)\leq U.
\end{equation}
The second chirality has been used only through $W_-\geq0$; no monopole
class in the opposite orientation has been assumed.

\subsection{Recovery by shrinking the surgery regions}

\begin{proposition}[Surgery recovery]\label{prop:vg-recovery}
On every fixed manifold \eqref{eq:vg-surgery-data} there are smooth metrics
$g_\varepsilon$, $\varepsilon\downarrow0$, for which
\begin{equation}\label{eq:vg-recovery-limits}
 \begin{gathered}
 V_\varepsilon\longrightarrow2\pi^2A,\qquad
 I_\varepsilon\longrightarrow-8\pi^2A,\qquad
 S_\varepsilon\longrightarrow32\pi^2A,\\
 W_{+,\varepsilon}\longrightarrow\frac{4\pi^2A}{3},\qquad
 R_\varepsilon\longrightarrow8\pi^2(A+Q),\qquad
 P_\varepsilon\longrightarrow8\pi^2Q.
 \end{gathered}
\end{equation}
Consequently
\begin{equation}\label{eq:vg-recovery-bound}
 \inf_g\E(g)\leq\frac{4h}{1+h}.
\end{equation}
Neither this inequality nor its construction asserts attainment.
\end{proposition}

\begin{proof}
We use the scalar-flat anti-self-dual insertions in the proof of
\cite[Theorem~4.3]{LeBrun2001}, and record the additional limits needed
for \eqref{eq:vg-defect-formulas}.  Equip
$\overline{\mathbb{CP}}^{\,2}$ with the oppositely oriented Fubini--Study
metric, and $S^1\times S^3$ with a round product metric.  These metrics
have positive scalar curvature and $W^+=0$.  Multiplication by the square
of a positive Green function of the conformal Laplacian at a point gives
a complete scalar-flat anti-self-dual metric on the punctured manifold.
In asymptotically Euclidean coordinates $y$ it has the expansion
\[
 h_{ab}(y)=\delta_{ab}+O(|y|^{-2}),\qquad
 \partial^m(h_{ab}-\delta_{ab})=O(|y|^{-2-m}),\quad m=1,2.
\]
These are the models used in the cited construction.

Multiply each model by $\varepsilon^4$ and use coordinates
$x=\varepsilon^2y$ on its end.  Its coefficients then satisfy
\begin{equation}\label{eq:vg-scaled-model}
 \partial^m(h_{\varepsilon,ab}-\delta_{ab})
       =O(\varepsilon^4|x|^{-2-m}),\qquad m=0,1,2.
\end{equation}
Choose $k+\ell$ distinct points in $X$ and disjoint normal-coordinate
balls about them.  At each point remove a ball of radius $\varepsilon$
and attach the appropriate model, truncated at its end.  On
$\varepsilon\leq r=|x|\leq2\varepsilon$, use
\[
 g_\varepsilon
   =f(r/\varepsilon)h_\varepsilon
       +(1-f(r/\varepsilon))g_X,
\]
where $f$ is a fixed smooth cutoff equal to one near the inner boundary
and zero near the outer boundary.  Away from these annuli the metric
is the model or $g_X$, respectively.  The asymptotic charts and the
orientation of the attachments give precisely the connected sum in
\eqref{eq:vg-surgery-data}.

In the normal coordinates, $g_X=\delta+O(r^2)$, with first derivative
$O(r)$ and second derivative $O(1)$.  Equation
\eqref{eq:vg-scaled-model} gives the same orders on $r\asymp\varepsilon$.
Differentiating the cutoff twice therefore shows that the interpolated
metric is uniformly equivalent to the Euclidean metric, has first
derivatives $O(\varepsilon)$ and second derivatives $O(1)$ on the
annulus.  Its curvature is consequently $O(1)$ there, and the annular
volume is $O(\varepsilon^4)$.  The constants here may depend on the
fixed manifold and the finite set of attachment points, which is all
that is required.

The inserted region also has volume $O(\varepsilon^4)$.  To see this,
the unscaled region extends to $|y|=O(\varepsilon^{-1})$ and has volume
$O(\varepsilon^{-4})$ by its Euclidean asymptotics; multiplication of a
four-dimensional metric by $\varepsilon^4$ multiplies volume by
$\varepsilon^8$.  The compact part of the model has still smaller
volume $O(\varepsilon^8)$.  On this entire inserted region $s=0$ and
$W^+=0$.  The exterior is unchanged and exhausts $X$ up to balls of
vanishing volume.  These observations prove the first four limits of
\eqref{eq:vg-recovery-limits}, including the signed integral $I$.

Solving \eqref{eq:vg-chiral-identities} for $P$ and $R$ gives
\[
 P=4W_++\frac S{12}-8\pi^2T_+,
 \qquad R=4W_++\frac S3-8\pi^2T_+.
\]
Together with $T_+=A-Q$, the first four limits give
$R_\varepsilon\to8\pi^2(A+Q)$ and
$P_\varepsilon\to8\pi^2Q$.  Finally,
\[
 \lim_{\varepsilon\downarrow0}\E(g_\varepsilon)
 =4-\frac{(8\pi^2A)^2}{(2\pi^2A)(8\pi^2(A+Q))}
 =\frac{4Q}{A+Q},
\]
which proves \eqref{eq:vg-recovery-bound}.
\end{proof}

\subsection{An integer-valued family in the two-chirality window}

Define
\begin{equation}\label{eq:vg-F}
 F(t)=\frac{(4-\sqrt{1-3t})^2}{9},\qquad 0\leq t<\frac13.
\end{equation}

\begin{theorem}[Balanced surgery family]\label{thm:vg-balanced-family}
There are finite unramified covers $X_j$ of a fixed closed
complex-hyperbolic surface, and nonnegative integers $k_j,\ell_j$, such
that, with the notation \eqref{eq:vg-surgery-data} on
\[
 M_j=X_j\#k_j\overline{\mathbb{CP}}^{\,2}\#\ell_j(S^1\times S^3),
\]
one has
\begin{equation}\label{eq:vg-balanced-window}
 \begin{gathered}
 h_j\longrightarrow0,\qquad Q_j\longrightarrow\infty,
 \qquad 0<U_j-F(h_j)=o(h_j^2),\\
 \frac{k_j}{Q_j}\longrightarrow\frac29,
 \qquad \frac{\ell_j}{Q_j}\longrightarrow\frac7{36}.
 \end{gathered}
\end{equation}
There is a constant $c_*>0$, independent of $j$, such that for all
sufficiently large $j$,
\begin{equation}\label{eq:vg-linear-sandwich}
 c_*h_j\leq\inf_g\J(g)\leq\inf_g\E(g)
                         \leq\frac{4h_j}{1+h_j}.
\end{equation}
In particular none of these sufficiently late $M_j$ admits an Einstein
metric, although their integral defects tend to zero.
\end{theorem}

\begin{proof}
Write the base surface as $X_0=\mathbb{CH}^2/\Gamma$.
The group $\Gamma$ is infinite and finitely generated, and its adjoint
representation is faithful and linear in characteristic zero.  It is
residually finite by Mal'cev's theorem \cite{Malcev1940}.
Its finite-index subgroups have unbounded indices: otherwise finite
generation would give only finitely many such subgroups, since each
subgroup of index $d$ arises from an action on $d$ points and there are
only finitely many homomorphisms from a finitely generated group to
$\mathfrak S_d$.  Their intersection would then have finite index.
Residual finiteness says that this intersection is trivial, contradicting
the infinitude of $\Gamma$.  Choose connected covers $X_j\to X_0$
with
\[
 A_j=c_1^2(X_j)=[\Gamma:\pi_1(X_j)]c_1^2(X_0)\geq j^4.
\]
No prescribed cover degree or divisibility condition is needed.

Put $\eta_j=1/j$, and for $t>0$ small define
\begin{equation}\label{eq:vg-rounding-choice}
 x(t)=\frac{F(t)-1+3t}{18},\qquad
 y(t)=\frac{15t-F(t)+1}{72},\qquad
 k_j=\lfloor A_jx(\eta_j)\rfloor,
 \quad \ell_j=\lfloor A_jy(\eta_j)\rfloor-2.
\end{equation}
Taylor expansion gives
\[
 F(t)=1+t+t^2+O(t^3),\qquad
 x(t)=\frac29t+O(t^2),\qquad
 y(t)=\frac7{36}t+O(t^2).
\]
Since $A_j\eta_j\geq j^3$, both integers are nonnegative for large $j$.
Let $a_j,b_j\in[0,1)$ be the fractional parts of $A_jx(\eta_j)$ and
$A_jy(\eta_j)$, respectively.  The exact identities
\[
 x(t)+4y(t)=t,\qquad 1+15x(t)-12y(t)=F(t)
\]
give
\begin{equation}\label{eq:vg-rounding-errors}
 h_j=\eta_j-\frac{a_j+4b_j+8}{A_j},\qquad
 U_j=F(\eta_j)+\frac{-15a_j+12b_j+24}{A_j}.
\end{equation}
For the second equality, $k_j/A_j\to0$, so
$\tau(M_j)=A_j/9-k_j>0$ for large $j$.  Hence the smaller quantity in
the definition of $U_j$ really is $T_{-,j}$.

For $0\leq t\leq19/300$,
\[
 0<F'(t)=\frac{4-\sqrt{1-3t}}{3\sqrt{1-3t}}<2.
\]
Also $0<\eta_j-h_j<13/A_j$ and
$9<-15a_j+12b_j+24<36$.  The mean value theorem applied to
\eqref{eq:vg-rounding-errors} therefore yields
\[
 0<U_j-F(h_j)<\frac{36+26}{A_j}
                    =\frac{62}{A_j}=o(h_j^2).
\]
Here $h_j\sim1/j$, because $A_j\geq j^4$.
The definitions also give $Q_j=A_jh_j\to\infty$ and the two ratios
in \eqref{eq:vg-balanced-window}.  In particular,
\begin{equation}\label{eq:vg-U-expansion}
 U_j=1+h_j+O(h_j^2).
\end{equation}

For every metric on $M_j$, the actual class
\eqref{eq:vg-actual-class} satisfies the hypotheses of
Theorem~\ref{thm:linear-window}.  Its linear estimate gives
\[
 \J(g)+(U_j-1-h_j)_+\geq c h_j.
\]
By \eqref{eq:vg-U-expansion}, the second term is $O(h_j^2)$, independently
of the metric.  After discarding finitely many indices this proves
$\J(g)\geq(c/2)h_j$ for every $g$.  Proposition~\ref{prop:vg-recovery}
and $\J\leq\E$ supply the other inequalities in
\eqref{eq:vg-linear-sandwich}.  An Einstein metric with nonzero Ricci
curvature has $\J=0$, contradicting the lower bound; Ricci-flatness
is separately excluded by \eqref{eq:vg-scalar-budget}.
\end{proof}

The recovery is performed on each fixed $M_j$.  Choosing its surgery
parameter after $j$ gives a diagonal sequence whose defect differs from
$4h_j/(1+h_j)$ by any prescribed $o(h_j)$ error.  No uniform lower
bound on the attachment radii is needed.  Conversely, these integral
estimates make no assertion that the pointwise Ricci error tends to zero.

\subsection{Conformal minimization and the traceless-Ricci problem}

The following comparison applies whenever every conformal class has an
actual positive monopole projection, and hence applies to all the surgery
manifolds above.

\begin{proposition}[Exact conformal comparison]\label{prop:vg-conformal-infimum}
Suppose that every conformal class on a closed four-manifold $M$ contains
an actual conformally solvable Spin$^c$ structure with positive self-dual
square.  Every conformal class then has a unique representative
$\widehat g$ with scalar curvature $-4$.  If
$\alpha=\inf_g\E(g)$, then
\begin{equation}\label{eq:vg-infima-exact}
 \inf_g\J(g)=\min\{1,\alpha\}.
\end{equation}
In particular, if $\alpha<1$, the infima of $\E$ and $\J$ coincide,
and both may be computed on the scalar-$-4$ representatives.
For $g=u^2\widehat g$, set $\delta=S(g)-S(\widehat g)$.  One has the
exact identities
\begin{equation}\label{eq:vg-conformal-scalar}
 \delta=48\int_Mu^{-2}|du|_{\widehat g}^2\,d\mu_{\widehat g}
       +36\int_M(u^{-1}\Delta_{\widehat g}u)^2\,d\mu_{\widehat g},
\end{equation}
\begin{equation}\label{eq:vg-conformal-P-R}
 P(g)=P(\widehat g)+\frac\delta{12},\qquad
 R(g)=R(\widehat g)+\frac\delta3,
\end{equation}
and, when $\alpha<1$,
\begin{equation}\label{eq:vg-conformal-excess}
 \E(g)-\alpha
 =\frac{R(\widehat g)}{R(g)}\bigl(\J(\widehat g)-\alpha\bigr)
     +(1-\alpha)\frac{\delta}{3R(g)}+\frac{L_s(g)}{R(g)}.
\end{equation}
All three terms on the right of \eqref{eq:vg-conformal-excess} are
nonnegative.
\end{proposition}

\begin{proof}
The Yamabe theorem \cite{LeeParker1987} gives a constant-scalar-curvature representative in
each conformal class.  Its scalar curvature cannot be nonnegative:
the unperturbed Weitzenb\"ock formula would force the spinor to vanish,
then its self-dual determinant curvature would vanish, contradicting
the positive self-dual square.  A constant rescaling therefore gives
$s_{\widehat g}=-4$.  For uniqueness, if both $\widehat g$ and
$u^2\widehat g$ have scalar curvature $-4$, then, with
$\Delta=d^*d$,
\begin{equation}\label{eq:vg-conformal-scalar-law}
 -4u^3=6\Delta_{\widehat g}u-4u.
\end{equation}
At a maximum of $u$ the positive Laplacian is nonnegative, so $u\leq1$;
at a minimum it is nonpositive, so $u\geq1$.  Thus $u=1$.

For a general $g=u^2\widehat g$, the transformation law is
$s_gu^3=6\Delta_{\widehat g}u-4u$.  Squaring it and using
$d\mu_g=u^4d\mu_{\widehat g}$ yields
\[
 S(g)=16\Vol(\widehat g)-48\int_Mu^{-1}\Delta_{\widehat g}u
                   +36\int_M(u^{-1}\Delta_{\widehat g}u)^2.
\]
Integration by parts gives
$\int u^{-1}\Delta u=-\int u^{-2}|du|^2$, proving
\eqref{eq:vg-conformal-scalar}.  In particular $\delta\geq0$, with
equality exactly when $u$ is constant.  The conformal invariance of
$W_+$ and \eqref{eq:vg-chiral-identities} give
$P(g)-P(\widehat g)=\delta/12$, and then $R=P+S/4$ gives the other
identity in \eqref{eq:vg-conformal-P-R}.  Consequently
\begin{equation}\label{eq:vg-J-convex}
 \J(g)=\frac{R(\widehat g)\J(\widehat g)+\delta/3}
                   {R(\widehat g)+\delta/3}.
\end{equation}
Thus $\J(g)$ lies between $\J(\widehat g)$ and one.  Since
$\E(\widehat g)=\J(\widehat g)$, whenever $\J(g)<1$ we obtain
\begin{equation}\label{eq:vg-Yamabe-improves}
 \E(\widehat g)=\J(\widehat g)\leq\J(g)\leq\E(g).
\end{equation}

We also need to see that the endpoint one in
\eqref{eq:vg-J-convex} can be approached.  Choose a coordinate ball,
a smooth cutoff $\chi$ supported in it and equal to one on a smaller
ball, and a fixed $0<\epsilon<1/2$.  The positive functions
$u_n=1+\epsilon\chi\sin(nx^1)$ are bounded above and below independently
of $n$.  On the smaller ball the leading term of
$\Delta_{\widehat g}u_n$ is
$\epsilon n^2\widehat g^{11}\sin(nx^1)$, and its remainder is $O(n)$
in $L^2$.  The integral of
$(\widehat g^{11})^2\sin^2(nx^1)$ on that ball tends to half the
integral of $(\widehat g^{11})^2$, by integration of the oscillatory
cosine term.  It is therefore bounded below by a positive constant for
large $n$.  Hence
$\int(u_n^{-1}\Delta u_n)^2\to\infty$ and \eqref{eq:vg-conformal-scalar}
gives $\delta_n\to\infty$.  Equation \eqref{eq:vg-J-convex} now shows
$\J(u_n^2\widehat g)\to1$.

It follows, first within a fixed class, that
\begin{equation}\label{eq:vg-class-infimum}
 \inf_{g\in[\widehat g]}\J(g)=\min\{\J(\widehat g),1\}.
\end{equation}
If $\J(\widehat g)<1$, the minimum is attained exactly by constant
rescalings of $\widehat g$; if $\J(\widehat g)=1$, every conformal metric
has $\J=1$; and if $\J(\widehat g)>1$, the infimum one is not attained.
Globally, $\J\leq\E$ and the preceding sequence give
$\inf\J\leq\min\{\alpha,1\}$.  On the other hand, a metric with
$\J(g)<\min\{\alpha,1\}$ would, by
\eqref{eq:vg-Yamabe-improves}, produce
$\E(\widehat g)<\alpha$, which is impossible.  This proves
\eqref{eq:vg-infima-exact} and the restriction to scalar-$-4$ metrics
when $\alpha<1$.

Finally, add $L_s/R$ to \eqref{eq:vg-J-convex}, subtract $\alpha$, and
use $R=R(\widehat g)+\delta/3$.  The result is
\eqref{eq:vg-conformal-excess}.  The first term is nonnegative because
$\J(\widehat g)=\E(\widehat g)\geq\alpha$, and the other two are
nonnegative by their definitions.
\end{proof}

In particular, for an $\E$-minimizing sequence on a fixed such manifold
with $\alpha<1$, \eqref{eq:vg-conformal-excess} gives
\[
 \frac{\delta}{R}\longrightarrow0,\qquad
 \frac{L_s}{R}\longrightarrow0,\qquad
 \frac{R(\widehat g)}R\longrightarrow1,
 \qquad \J(\widehat g)\longrightarrow\alpha.
\]
Its late terms may be replaced by their scalar-$-4$ representatives
without increasing the defect.  These conclusions control discarded
energies; they do not assert convergence of the conformal factors or
compactness of the resulting metrics.

\begin{theorem}[Reduction to traceless and full curvature energy]
\label{thm:vg-first-order-reduction}
For the family in Theorem~\ref{thm:vg-balanced-family}, put
\begin{equation}\label{eq:vg-alpha-rho}
 \alpha_j=\inf_{g\text{ on }M_j}\E(g),\qquad
 \rho_j=\frac{1}{8\pi^2A_j}
                    \inf_{g\text{ on }M_j}\int_{M_j}|\Ric_g^\circ|^2.
\end{equation}
For all sufficiently large $j$,
\begin{equation}\label{eq:vg-rho-bounds}
 \inf_g\J(g)=\alpha_j,\qquad
 \frac{4\rho_j}{U_j+4\rho_j}
                  \leq\alpha_j\leq\frac{4\rho_j}{1+\rho_j}.
\end{equation}
Consequently
\begin{equation}\label{eq:vg-alpha-four-rho}
 \rho_j=O(h_j),\qquad \alpha_j=4\rho_j+O(h_j^2).
\end{equation}
If $\Rm$ denotes the curvature operator on $\Lambda^2$, with its
Hilbert--Schmidt norm, then
\begin{equation}\label{eq:vg-full-curvature-infimum}
 \inf_g\int_{M_j}|\Rm_g|^2
      =8\pi^2\chi(M_j)+8\pi^2A_j\rho_j.
\end{equation}
In particular, the assertion $\alpha_j/h_j\to4$ is equivalent to either
of the assertions
\begin{equation}\label{eq:vg-open-equivalences}
 \inf_g\int_{M_j}|\Ric_g^\circ|^2
          =8\pi^2Q_j+o(Q_j),\qquad
 \inf_g\int_{M_j}|\Rm_g|^2
          =8\pi^2\chi(M_j)+8\pi^2Q_j+o(Q_j).
\end{equation}
\end{theorem}

\begin{proof}
The recovery bound makes $\alpha_j<1$ for large $j$, so
Proposition~\ref{prop:vg-conformal-infimum} gives the first equality in
\eqref{eq:vg-rho-bounds}.  Fix such a $j$, temporarily suppress its
index, and put $x=P/(8\pi^2A)$.  Equation \eqref{eq:vg-two-sided-R}
implies
\[
 \frac{4x}{U+4x}\leq\J(g)\leq\frac{4x}{1+x}.
\]
Both functions of $x$ are continuous and increasing on $[0,\infty)$,
since $U>0$.  Their infima over the metric values of $x$ are therefore
their values at $\rho$.  More explicitly, $x\geq\rho$ gives the lower
bound, and a sequence with $x\downarrow\rho$ gives the upper bound
after taking the infimum of $\J$.  No common minimizing metric for
$P$ and $\J$ is required.  This proves \eqref{eq:vg-rho-bounds}.

Its lower inequality and $\alpha<1$ give
\[
 \rho\leq\frac{\alpha U}{4(1-\alpha)}.
\]
Since $\alpha=O(h)$ and $U=1+O(h)$ by
\eqref{eq:vg-linear-sandwich} and \eqref{eq:vg-U-expansion}, this proves
$\rho=O(h)$.  Both outer terms of \eqref{eq:vg-rho-bounds} are now
$4\rho+O(h^2)$, which proves \eqref{eq:vg-alpha-four-rho}.

The orthogonal decomposition of the curvature operator and the
Chern--Gauss--Bonnet formula give, respectively,
\[
 \int_M|\Rm_g|^2=W_++W_-+\frac S{24}+\frac P2,
 \qquad
 8\pi^2\chi(M)=W_++W_-+\frac S{24}-\frac P2.
\]
Subtracting gives the exact identity
$\int_M|\Rm_g|^2=8\pi^2\chi(M)+P$ for every metric, and hence
\eqref{eq:vg-full-curvature-infimum}.  If curvature is instead regarded
as a covariant four-tensor with the full tensor norm, its squared norm
is four times the operator norm used here; the constants in
\eqref{eq:vg-full-curvature-infimum} must then also be multiplied by four.
Finally divide \eqref{eq:vg-alpha-four-rho} by $h_j$, use
$A_jh_j=Q_j$, and then apply \eqref{eq:vg-full-curvature-infimum}.
This proves the two equivalences in \eqref{eq:vg-open-equivalences}.
\end{proof}

\subsection{An exact subfamily and the unresolved balanced coefficient}

The upper coefficient can be determined when handles dominate the surgery.
This fact does not determine it for the balanced family.

\begin{proposition}[A scalar--Weyl identity]\label{prop:vg-handle-identity}
For the manifold \eqref{eq:vg-surgery-data}, set
\[
 \theta=\frac{4\ell-5k}{A}=h-\frac{6k}{A},\qquad
 S_0=32\pi^2A,\qquad B=\int_M|W^-_g|^2.
\]
If $0\leq\theta\leq1/3$, every smooth metric satisfies
\begin{equation}\label{eq:vg-handle-identity}
 \E(g)-\frac{4\theta}{1+\theta}
 =\frac{(1-3\theta)(S-S_0)/3+16B+(1+\theta)L_s}
                   {(1+\theta)R}.
\end{equation}
In particular, for all values of $k,\ell$,
\begin{equation}\label{eq:vg-piecewise-bound}
 \E(g)\geq
 \begin{cases}
 0,&\theta\leq0,\\
 4\theta/(1+\theta),&0\leq\theta\leq1/3,\\
 1,&\theta\geq1/3.
 \end{cases}
\end{equation}
\end{proposition}

\begin{proof}
The signature formula and \eqref{eq:vg-chiral-identities} give
\[
 W_+=B+12\pi^2(A/9-k),\qquad
 P=4W_++S/12-8\pi^2(A-k-4\ell).
\]
Thus
\begin{equation}\label{eq:vg-P-theta}
 P=\frac S{12}+4B+8\pi^2A\left(\theta-\frac13\right).
\end{equation}
By \eqref{eq:vg-defect-formulas}, the numerator of
$(1+\theta)R\bigl(\E-4\theta/(1+\theta)\bigr)$ is
\[
 (1+\theta)(4P+L_s)-4\theta(P+S/4)
 =4P-\theta S+(1+\theta)L_s.
\]
Substitution of \eqref{eq:vg-P-theta} turns this into
$(1-3\theta)(S-S_0)/3+16B+(1+\theta)L_s$, proving
\eqref{eq:vg-handle-identity}.  Every term is nonnegative when
$0\leq\theta\leq1/3$, by \eqref{eq:vg-scalar-budget}.  For
$\theta\geq1/3$, the same substitution gives instead
\[
 \E-1=\frac{3P-S/4+L_s}{R}
       =\frac{8\pi^2A(3\theta-1)+12B+L_s}{R}\geq0.
\]
For $\theta\leq0$ the definition of $\E$ gives the stated bound.
\end{proof}

\begin{theorem}[Exact handle-only infimum]\label{thm:vg-handle-sharp}
If $k=0$ and $0\leq h=4\ell/A\leq1/3$, then
\begin{equation}\label{eq:vg-handle-sharp}
 \inf_g\E(g)=\inf_g\J(g)=\frac{4h}{1+h}.
\end{equation}
For $h=0$, the complex-hyperbolic metric attains the infima.
For $0<h<1/3$, neither infimum is attained by a smooth metric.
The argument below makes no assertion of attainment or nonattainment
at $h=1/3$.
\end{theorem}

\begin{proof}
Here $\theta=h$, so \eqref{eq:vg-handle-identity} is an all-metric lower
bound $\E\geq4h/(1+h)$.  Proposition~\ref{prop:vg-recovery} gives
the matching upper bound on the same smooth manifold.  Since its value
is at most one, \eqref{eq:vg-infima-exact} gives the assertion about
$\inf\J$ as well.  If $h=0$, $M=X$ and $g_X$ has zero defect.

Suppose now that $0<h<1/3$ and equality is attained for $\E$.
The coefficients in \eqref{eq:vg-handle-identity} force
$S=S_0$, $B=0$ and $L_s=0$.  Take the actual ordinary monopole
structure supplied above; since there are no exceptional summands,
its determinant class $a$ has $a^2=A$.  For a solution $(\mathcal A,\Phi)$,
the inequalities used to prove \eqref{eq:vg-scalar-budget} become
\[
 S_0\leq32\pi^2p\leq H_\Phi\leq S=S_0.
\]
Every inequality is therefore an equality.  The Weitzenb\"ock formula
then forces $\int|\nabla_{\mathcal A}\Phi|^2=0$.  Hence $\Phi$ is a
nonzero parallel positive spinor and is nowhere zero.

For the positive spinor bundle of a Spin$^c$ structure with determinant
class $a$, the characteristic-class identity is
\[
 4c_2(S_a^+)[M]=a^2-(2\chi(M)+3\tau(M)).
\]
For example, it follows directly from
$\mathfrak{su}(S_a^+)\cong\Lambda^+$, together with
$p_1(\mathfrak{su}(E))=c_1(E)^2-4c_2(E)$ for a complex rank-two bundle
and $p_1(\Lambda^+)[M]=2\chi(M)+3\tau(M)$.
Here it gives $c_2(S_a^+)[M]=\ell>0$.  A nowhere-zero section of a
complex rank-two bundle spans a trivial line subbundle, so its second
Chern class vanishes, a contradiction.  Thus $\E$ cannot attain its
infimum.  If $\J$ attained that same value, which is less than one,
\eqref{eq:vg-Yamabe-improves} would produce an $\E$-minimizer, again a
contradiction.  At $h=1/3$ the coefficient of $S-S_0$ vanishes, so the
preceding scalar-equality argument is not available.
\end{proof}

More generally, if $h_j\downarrow0$ and $k_j=o(Q_j)$ in any such surgery
sequence, then $\theta_j/h_j=1-6k_j/Q_j\to1$.
For large $j$, $0<\theta_j<1/3$, and
\eqref{eq:vg-piecewise-bound} and \eqref{eq:vg-recovery-bound} give
\begin{equation}\label{eq:vg-handle-dominant}
 \frac{4\theta_j}{h_j(1+\theta_j)}
   \leq\frac{\inf_g\E(g)}{h_j}\leq\frac4{1+h_j}.
\end{equation}
Both outer terms tend to four.  Thus the first-order coefficient is four
on these handle-dominant sequences.

For the balanced family, in contrast,
$k_j/Q_j\to2/9$, so $\theta_j/h_j\to-1/3$.
The elementary estimate \eqref{eq:vg-piecewise-bound} is then only
the trivial lower bound zero.  Theorem~\ref{thm:linear-window} supplies
the positive linear order in \eqref{eq:vg-linear-sandwich}, but it does
not determine the best coefficient.  The remaining variational question
is to determine the asymptotic value, or the distinct lower and upper
limits if no limit exists, of
\begin{equation}\label{eq:vg-open-coefficient}
 \frac{\alpha_j}{h_j}
 =\frac1{h_j}\inf_{g\text{ on }M_j}\E(g)
 =\frac1{h_j}\inf_{g\text{ on }M_j}\J(g).
\end{equation}
What is proved is
\[
 0<c_*\leq\liminf_{j\to\infty}\frac{\alpha_j}{h_j}
     \leq\limsup_{j\to\infty}\frac{\alpha_j}{h_j}\leq4.
\]
The assertion that both limits are four is precisely the all-metric
energy assertion \eqref{eq:vg-open-equivalences}; it is not a consequence
of the recovery sequence.  Conversely, an upper construction with a
smaller coefficient would require a matching lower analysis to determine
the new optimum.  A classification of all degenerations or a full
limiting variational theory would be stronger conclusions, not premises
of the equivalences proved here.
